%%%%%%%%%%  scan idx 608  |  folio 595  %%%%%%%%%%
\section*{\underline{EXPOSÉ XII}\\[3pt] \underline{\underline{UN THÉORÈME DE REPRÉSENTABILITÉ RELATIVE SUR}}\\[3pt] \underline{\underline{LE FONCTEUR DE PICARD}}}
\begin{center}
\underline{par M. Raynaud (rédigé par S. Kleiman)}
\end{center}

\subsection*{Avertissement}

Le but des deux exposés XII, XIII est de donner des démonstrations des neuf théorèmes de finitude pour le foncteur de Picard énoncés par A. Grothendieck dans les commentaires (p. C 07-C 11) de ses «~Fondements de la Géométrie Algébrique~» (cité FGA). Dans le premier exposé, on démontre le premier théorème de finitude sous la forme suivante plus précise que celle énoncée dans FGA~: soient \(X\) et \(Y\) deux schémas propres sur une base noethérienne intègre \(S\), \(f:X\longrightarrow Y\) un \(S\)-morphisme surjectif. Alors il existe un ouvert non vide \(V\) de \(S\) au-dessus duquel le morphisme induit \(f^*:\Pic_{Y/S}\longrightarrow \Pic_{X/S}\) est représentable par un morphisme quasi-affine de présentation finie. Dans le deuxième exposé, on démontre le reste.

Les deux exposés diffèrent notablement par leur technique de démonstration. Dans le premier, on emploie de façon essentielle le dévissage de Oort, la descente non plate et la représentation des foncteurs non-ramifiés. Dans le deuxième, on se réduit d'abord au cas où les schémas \(X/S\) qui interviennent sont projectifs et plats à fibres géométriquement intègres, ou même au besoin, normales ou lisses, toujours en appliquant le premier théorème de finitude à un morphisme \(X'\to X\) convenable. Ensuite on emploie l'outil des \((b)\)-faisceaux forgé par I. Castelnuovo et aiguisé et utilisé par T. Matsusaka, D. Mumford, Y. Nakai et S. Kleiman. Cet outil permet notamment de se débarrasser des coordonnées de Chow, du morphisme de Frobenius et des théorèmes de Lefschetz, utilisés dans les démonstrations initiales de Grothendieck.

Le présent exposé XII a été présenté oralement d'après Grothendieck par M. Raynaud, que le rédacteur tient à remercier pour lui avoir communiqué ses notes personnelles.

